\chapter{Randomized Hashing and Streaming}
\label{ch:stage4-hashing-streaming}

This supplement connects the book's universal hashing, perfect hashing,
fingerprinting, and streaming sections to the data structures used in modern
systems.  Each structure has a different contract: membership, frequency,
heavy hitters, or distinct counting.  Students should not treat these
contracts as interchangeable.

\section{Cuckoo hashing}

Choose two hash functions $h_1,h_2$ and store each key $x$ in one of the two
locations $h_1(x)$ or $h_2(x)$.  An insertion may evict the key currently in
one location and move it to its alternate location.  If a cycle is detected,
rehash with fresh functions.

The standard random-graph representation has one vertex per table position
and one edge $(h_1(x),h_2(x))$ per key.  Successful placement corresponds to
orienting every edge toward a distinct endpoint.  A component with more
edges than vertices prevents such an orientation.  Below the load threshold,
the random graph has no complex obstruction with high probability, while
above the threshold insertion failures become likely.  The exact threshold
depends on the table layout and load convention, so implementations must
state their resize policy.

\begin{proposition}[Cuckoo lookup]
If the table is maintained with two candidate positions, a lookup examines
at most two locations and therefore takes $O(1)$ worst-case time.  Insertions
have expected $O(1)$ displacement work when the load is kept below a fixed
threshold and failed insertions trigger rehashing.
\end{proposition}

\section{Consistent hashing}

Place both servers and keys on a circular hash space.  A key is assigned to
the first server encountered clockwise.  Adding or removing one server then
changes only the keys in its predecessor interval, rather than remapping all
keys.

With $m$ independently hashed servers and $n$ independently hashed keys, the
expected fraction of keys moved when one server is added is $1/(m+1)$.
Virtual replicas reduce load imbalance: with $r$ replicas per server, the
largest interval decreases as $r$ grows, but the replicas are not magically
independent if the hash construction correlates them.

\section{Bloom filters}

A Bloom filter uses an $m$-bit array and $k$ independent hash functions.  To
insert $x$, set the $k$ indexed bits.  A query returns ``absent'' if any of
the bits is zero and ``possibly present'' otherwise.

There are no false negatives unless deletion is performed incorrectly.  If
$n$ items have been inserted, the false-positive probability is
\[
 \left(1-\left(1-\frac1m\right)^{kn}\right)^k
 \approx (1-e^{-kn/m})^k.
\]
The approximation is minimized near $k=(m/n)\ln 2$, giving roughly
$(0.6185)^{m/n}$.  Standard Bloom filters do not support deletion; counting
Bloom filters or cuckoo filters change the space and error analysis.

\section{Count-Min Sketch}

For a turnstile or insertion-only stream, choose $d$ pairwise-independent
row hashes into $w$ counters.  On update $(x,\Delta)$, increment counter
$(i,h_i(x))$ by $\Delta$ in every row.  Estimate a point frequency by the
minimum of its row counters.

\begin{theorem}[Count-Min guarantee]
For a nonnegative insertion-only stream with total mass $F_1$, choosing
$w=\lceil e/\varepsilon\rceil$ and
$d=\lceil\ln(1/\delta)\rceil$ gives, for every fixed query $x$,
\[
 f_x\leq\widehat f_x\leq f_x+\varepsilon F_1
\]
with probability at least $1-\delta$.
\end{theorem}

\begin{proof}
In one row, the expected collision mass added to $x$ is at most
$(F_1-f_x)/w\leq F_1/w$.  Markov's inequality makes the probability that
this noise exceeds $\varepsilon F_1$ at most $1/e$.  Taking the minimum of
$d$ independent rows reduces the failure probability to at most
$e^{-d}\leq\delta$.
\end{proof}

The same one-sided guarantee does not hold for arbitrary negative updates;
the turnstile model needs a different analysis or a signed sketch.

\section{Count-Sketch and signed updates}

Count-Sketch uses a bucket hash $h_i$ and an independent sign hash
$s_i(x)\in\{-1,+1\}$.  The update adds $s_i(x)\Delta$ to bucket $h_i(x)$,
and the estimate is the median of $s_i(x)$ times the queried buckets.
The sign cancellation makes the estimator unbiased under pairwise
independence.  With suitable width and depth, the error is controlled by the
$\ell_2$ mass of the frequency vector rather than its $\ell_1$ mass.

\section{Heavy hitters}

An item is an $\alpha$-heavy hitter if $f_x\geq\alpha F_1$.  The Misra--Gries
algorithm stores $k$ candidate counters and repeatedly decrements all
counters when the table is full.  At termination, every item with frequency
greater than $F_1/k$ is present among the candidates; a second pass gives
exact frequencies for those candidates.

Sketches can identify heavy hitters without a second pass, but the query
guarantee, update model, and space bound must be stated.  A common mistake is
to claim exact heavy-hitter identification from a Count-Min Sketch without
allowing for its additive error.

\section{Frequency moments}

For frequencies $(f_x)$, the $p$th frequency moment is
$F_p=\sum_x f_x^p$.  $F_1$ is stream length, $F_2$ measures concentration,
and $F_0$ is the number of distinct items.  The choice of estimator depends
strongly on whether the stream is insertion-only or allows deletions.

\section{Distinct-element estimation}

Flajolet--Martin stores the maximum trailing-zero statistic of hashed stream
items.  Repeating with independent hashes and aggregating robustly reduces
variance.  HyperLogLog refines this idea by keeping many registers and using
their harmonic-style aggregate.  The practical lesson is that a single
register is a noisy estimator; claims about relative error must specify the
number of registers and the aggregation rule.

\section{Network coding and randomized polling}

In randomized network coding, intermediate nodes transmit random linear
combinations of packets over a finite field.  A sink recovers the original
packets when a resulting transfer matrix has full rank.  The determinant of a
nonzero polynomial is nonzero under a random assignment with high probability
by Schwartz--Zippel, provided the field is sufficiently large.

In polling, sample stream positions or records using a specified sampling
distribution, then use inverse-probability weighting when estimating totals.
Uniform sampling without replacement creates dependence; variance estimates
must use the sampling design rather than incorrectly assuming independent
observations.

\section{Student checklist}

For each streaming structure, students should record:

\begin{enumerate}
\item insertion-only or turnstile update model;
\item query type and whether the guarantee is one-sided;
\item hash-family independence required;
\item space, update time, query time, and failure probability;
\item whether a second pass or rehashing is required;
\item the adversarial model for stream order and queried items.
\end{enumerate}

\begin{problem}[Bloom-filter optimization]
For a target false-positive probability $\delta$, derive the bits-per-item
needed by a Bloom filter using the optimal number of hash functions.
\end{problem}

\begin{problem}[Count-Min proof]
Prove the Count-Min guarantee for a fixed query and explain why a union bound
is required to guarantee all queries simultaneously.
\end{problem}

\begin{problem}[Cuckoo obstruction]
Show that a component with more edges than vertices in the cuckoo graph
prevents all keys from being assigned to distinct candidate locations.
\end{problem}

\begin{problem}[Streaming implementation]
Implement Count-Min, Count-Sketch, and Misra--Gries.  Compare their errors on
Zipfian, uniform, and adversarial streams.
\end{problem}
